% intro.tex — Introduction.

\section{Introduction}
\label{sec:intro}

When sensory cues signal that some spatial locations are more rewarding
than others, how should an observer allocate limited processing
resources? This question is central both to understanding biological
attention and to designing the experiments that probe it. An observer
informed that a change at one location is worth several times the
reward of a change elsewhere could, in principle, exploit that
information in more than one way, and the mechanisms differ in what they
cost.

The first mechanism is \emph{criterion adjustment}: shifting the
decision threshold at the high-value location to become more liberal,
accepting more false alarms in exchange for more hits
\cite{MullerFindlay1987, posner1980_orienting}. Criterion adjustment is
``free'' in the sense that it requires no perceptual reallocation --- the
observer simply changes how much evidence is needed to declare a change
at each location, and a shift at one location leaves sensitivity
elsewhere untouched.

The second mechanism is \emph{value-directed attention} (VDA):
re-allocating spatial attention toward the high-value location,
improving perceptual sensitivity ($\dprime$) there at the cost of
reduced sensitivity elsewhere \cite{CohenMaunsell2009}. Unlike criterion
adjustment, VDA involves a genuine perceptual trade-off --- resources
gained at one location must be taken from another. Reward-associated
stimuli can capture spatial attention
\cite{failing_theeuwes2018_selection_history, hickey2010_reward_salience_acc},
providing a candidate substrate for such re-allocation.

A third lever becomes available once we abandon the assumption that
detection noise is independent across locations. When the population
code carries noise that is correlated across locations, an observer can
adapt to value by \emph{decorrelating} the encoded signal --- and
attention is known to reduce interneuronal correlations in visual cortex
\cite{CohenMaunsell2009, RuffCohen2016, Srinath2021}. We therefore treat
the cross-location correlation $\corr$ as a model parameter rather than
fixing it at zero, which yields a three-lever decomposition of the
value response --- criterion, sensitivity, and decorrelation
(Section~\ref{sec:model}). The familiar criterion-versus-sensitivity
account is recovered as the $\corr = 0$ special case.

A further consideration is that the benefit and cost of attention are
not necessarily symmetric. Attentional enhancement at attended locations
and suppression at unattended locations are dissociable mechanisms that
can be modulated independently
\cite{reynolds_heeger2009_normalization, treue_martinez_trujillo1999_feature_attention, mcadams_maunsell1999_v4_tuning},
so a normative model that assumes symmetric benefit and cost may reach
conclusions specific to that assumption rather than general properties
of the task. We therefore equip the model with an explicit asymmetry
ratio $\Rsens > 0$ that parameterises the relative efficacy of
enhancement versus suppression. The normative question is then: under
what conditions does VDA gain meaningfully over an observer that has
already optimised its criteria and its validity-driven attention? If
the criterion lever captures most of the available reward, the neural
machinery supporting value-driven re-allocation is, for this class of
tasks, redundant, and experiments designed to observe VDA in that regime
will tend to produce null results. Conversely, identifying the
parameter regimes in which VDA is essential provides a recipe for
experimental designs that can isolate value-driven attentional effects
from criterion-based strategies.

Our analysis yields four conclusions
(Section~\ref{sec:results}). First, criterion adjustment is
\emph{typically} the dominant lever for encoding value: across a broad
parameter sweep the share of the reward gain it captures has a median
near $0.76$ and is concentrated in $[0.30, 1.00]$, a central tendency
with a substantial tail toward attention-dominated outcomes that grows
as correlated noise ($\corr > 0$) is admitted. Second, the advantage of VDA is \emph{non-monotonic} in the
benefit/cost ratio $\Rsens$, peaking in a cost-dominant regime where
value-blind attention is too expensive to deploy; we strengthen this
finding with a closed-form threshold $\rdagger(\val)$ that marks the
lower edge of the VDA-active band. Third, this advantage is concentrated
in a \emph{graded} regime --- low cue validity, high value contrast,
moderate-to-cost-dominant asymmetry, and low baseline sensitivity ---
which we map as an iso-VDA contour band over $(\valid, \val, \Rsens)$. Fourth, re-allocating
attention \emph{away} from the high-value location is not optimal under
predictive cues, a result that holds as a conditional theorem with the
explicit condition $\valid \ge 1/\Nloc$; outside that condition, in the
counter-predictive ``anti-cue'' regime ($\valid < 1/\Nloc$), inverted
allocation can become optimal, which we report as a new falsifiable
prediction of the model. Outside the graded VDA regime --- which
includes the operating point of many standard spatial-cueing paradigms
--- an observer that adjusts criteria optimally loses little by treating
attention as value-blind.
